%% ═══════════════════════════════════════════════════════════════════════
%%  Literature Survey
%%  Title: Channel Aging in Cell-Free Massive MIMO
%%  Font : Computer Modern (= CMU Serif — identical typeface)
%%  Colour: Black only
%%  Compiler: pdflatex
%% ═══════════════════════════════════════════════════════════════════════
\documentclass[12pt,a4paper]{article}

%% ── Font & Encoding ──────────────────────────────────────────────────
\usepackage[T1]{fontenc}
\usepackage[utf8]{inputenc}
%% Computer Modern is the same typeface family as CMU Serif (CMU Serif is
%% the OpenType packaging of Computer Modern). The output is visually identical.

%% ── Page Layout ──────────────────────────────────────────────────────
\usepackage[top=2.5cm,bottom=2.5cm,left=2.8cm,right=2.8cm]{geometry}

%% ── Math ─────────────────────────────────────────────────────────────
\usepackage{amsmath,amssymb}

%% ── Graphics & Figures ───────────────────────────────────────────────
\usepackage{graphicx}
\usepackage{caption}
\usepackage{subcaption}

%% ── Tables ───────────────────────────────────────────────────────────
\usepackage{booktabs}
\usepackage{array}
\usepackage{longtable}
\usepackage{multirow}

%% ── Typography & Layout ──────────────────────────────────────────────
\usepackage{setspace}

\usepackage{parskip}
\usepackage{enumitem}

%% ── Section Formatting ───────────────────────────────────────────────
\usepackage{titlesec}
\titleformat{\section}{\large\bfseries}{\thesection.}{0.6em}{}[\titlerule]
\titleformat{\subsection}{\normalsize\bfseries}{\thesubsection.}{0.5em}{}
\titleformat{\subsubsection}{\normalsize\itshape}{\thesubsubsection.}{0.5em}{}

%% ── References & Links ───────────────────────────────────────────────
\usepackage[hidelinks]{hyperref}
\usepackage{cite}

%% ── Spacing ──────────────────────────────────────────────────────────
\onehalfspacing

%% ── Math Macros ──────────────────────────────────────────────────────
\newcommand{\bH}{\mathbf{H}}
\newcommand{\bh}{\mathbf{h}}
\newcommand{\CC}{\mathbb{C}}
\newcommand{\bF}{\mathbf{F}}
\newcommand{\bK}{\mathbf{K}}
\newcommand{\bR}{\mathbf{R}}
\newcommand{\bY}{\mathbf{Y}}
\newcommand{\bPhi}{\boldsymbol{\Phi}}
\newcommand{\bI}{\mathbf{I}}
\newcommand{\bA}{\mathbf{A}}

%% ═══════════════════════════════════════════════════════════════════════
\begin{document}

%% ── Title Page ────────────────────────────────────────────────────────
\begin{titlepage}
  \centering
  \vspace*{2cm}
  {\LARGE\bfseries Literature Survey}\\[0.6em]
  \rule{\linewidth}{0.8pt}\\[0.8em]
  {\Large\bfseries Channel Aging in Cell-Free Massive MIMO:\\[0.5em]
   Towards Lightweight, Predictive CNN-Based\\[0.5em]
   Channel State Estimation for 6G Networks}\\[0.8em]
  \rule{\linewidth}{0.8pt}\\[2.0cm]

  {\normalsize\itshape
  Research Domain: AI-Native Physical Layer Design\\[0.3em]
  Research Area: 6G Wireless Communications}\\[2.5cm]

  \begin{minipage}{0.85\linewidth}
    \centering
    {\small\textbf{Refined Problem Statement:}}\\[0.5em]
    {\small\itshape
    ``Multi-step forward channel state prediction in distributed
    cell-free massive MIMO networks using computationally efficient
    CNNs designed for spatially irregular access point topologies,
    targeting real-time channel aging mitigation under high user
    mobility.''}
  \end{minipage}

  \vfill
  {\small\today}
\end{titlepage}

%% ── Abstract ──────────────────────────────────────────────────────────
\begin{abstract}
\noindent
This literature survey examines the problem of channel aging in
cell-free massive multiple-input multiple-output (MIMO) systems and
surveys existing approaches for mitigating it through channel prediction.
Beginning with foundational work on massive MIMO channel estimation,
the survey traces the evolution from classical model-based predictors
--- including autoregressive (AR) models and Kalman filters --- through
early machine learning approaches, to state-of-the-art deep learning
architectures employing convolutional neural networks (CNNs), long
short-term memory (LSTM) networks, and Transformer models. Special
attention is given to the cell-free architecture and the unique
challenges it introduces: irregular access point (AP) topology,
distributed processing constraints, and the need for computationally
efficient edge-deployable models. The survey synthesises findings from
over twenty peer-reviewed publications and identifies the critical
research gap that the proposed project seeks to address: the absence
of lightweight, multi-step predictive CNN architectures specifically
designed for the spatially irregular, distributed topology of cell-free
massive MIMO systems under high user mobility.
\end{abstract}

\newpage
\tableofcontents
\newpage

%% ═══════════════════════════════════════════════════════════════════════
\section{Introduction}

The demand for wireless connectivity has grown exponentially over the
past decade, driven by smartphones, IoT devices, autonomous vehicles,
and immersive extended reality applications. Fifth-generation (5G)
networks, while a substantial leap forward, are insufficient to meet
the projected traffic densities and latency requirements of the 2030
horizon. Sixth-generation (6G) networks are therefore under active
research, targeting terabit-per-second peak data rates, sub-millisecond
latency, and near-universal seamless coverage~\cite{Wang2024tutorial}.

Two of the most promising enabling technologies for 6G are massive
multiple-input multiple-output (mMIMO) systems and the cell-free (CF)
network architecture. Massive MIMO exploits large antenna arrays
--- typically 64 to 256 antennas per base station --- to serve many
users simultaneously through spatial multiplexing and
beamforming~\cite{Lu2014overview}. The cell-free paradigm extends this
concept by distributing a large number of access points (APs) across a
wide geographical area, all connected to a centralised processing unit
(CPU), collectively serving users without cell
boundaries~\cite{Ngo2017cellfree, Elhoushy2022survey}.

\begin{figure}[htbp]
  \centering
  \includegraphics[width=0.90\linewidth]{figures/fig_channel_aging.pdf}
  \caption{The channel aging problem: the channel state measured at
  pilot time ($\tilde{\bH}$) becomes outdated by the time precoding
  is applied, leading to beam misalignment and throughput loss.}
  \label{fig:aging}
\end{figure}

A fundamental operational requirement for both massive MIMO and
cell-free systems is accurate, real-time knowledge of the channel state
information (CSI) at the transmitter. CSI describes the complex
propagation characteristics between each transmit antenna and each user
device, encoding how the signal is attenuated, delayed, and
phase-shifted as it travels through the wireless environment. In
time-division duplex (TDD) systems, the transmitter estimates the
downlink channel by exploiting uplink-downlink reciprocity: the user
device transmits known pilot sequences, and the AP measures how these
pilots were distorted to reconstruct the channel matrix
$\bH \in \CC^{N \times K}$, where $N$ is the number of antennas and
$K$ is the number of subcarriers.

The critical challenge, illustrated in Figure~\ref{fig:aging}, is that
the process of pilot transmission, channel estimation, and precoder
computation takes several milliseconds. During this processing delay,
the physical environment changes --- users move, doors open, vehicles
pass --- causing the estimated channel $\tilde{\bH}$ to become stale
relative to the true channel at the time of data transmission. This
phenomenon is known as \textit{channel aging}, and it directly degrades
beamforming performance, spectral efficiency, and system
reliability~\cite{Zheng2021impact}.

Channel aging is particularly pronounced in cell-free systems for two
compounding reasons~\cite{Freitas2025ekf}: (1) centralised precoding
requires all AP-to-user channel matrices to be gathered at the CPU,
introducing additional fronthaul delays on top of the estimation delay;
and (2) the distributed architecture means that each user is tracked
by multiple independently aging channel links, all of which must remain
coherent for joint transmission to succeed.

This survey reviews the body of literature addressing channel aging
mitigation through predictive channel state estimation, with the goal
of establishing the research context and identifying the gap that
motivates the proposed project.

%% ═══════════════════════════════════════════════════════════════════════
\section{Massive MIMO: Foundations and Channel Estimation}

\subsection{Overview of Massive MIMO}

The concept of massive MIMO was formally popularised by Marzetta (2010)
and subsequently analysed in depth by Lu \textit{et al.}
\cite{Lu2014overview}. The fundamental insight is that as the number of
base station antennas $N$ grows large while the number of served users
$K$ remains fixed, inter-user interference averages out and the system
approaches a regime of near-orthogonal user channels. This allows simple
linear precoders such as maximum ratio transmission (MRT) to approach
optimal performance, substantially reducing computational
requirements~\cite{Lu2014overview}.

By the time of 5G standardisation, massive MIMO with 64-element active
antenna units had been commercially deployed in sub-6\,GHz bands.
Looking towards 6G, the concept has been extended to extremely
large-scale MIMO (XL-MIMO) with hundreds to thousands of antennas,
as surveyed in~\cite{Wang2024tutorial}, introducing additional
challenges in near-field propagation and spatially non-stationary
channel behaviour.

\subsection{Classical Channel Estimation Methods}

Channel estimation in massive MIMO systems relies on pilot sequences:
the user device transmits a known signal of length $\tau_p$ symbols,
and the receiver computes the channel estimate from the received pilot
observations. Two classical estimators are widely used:

\subsubsection{Least Squares (LS) Estimation}
The LS estimator computes:
\begin{equation}
  \hat{\bH}^{\mathrm{LS}} = \bY_p \bPhi^{\dagger}
\end{equation}
where $\bY_p$ is the received pilot matrix and $\bPhi$ is the pilot
sequence matrix. The LS estimator is unbiased and computationally
efficient but noise-amplifying.

\subsubsection{MMSE Estimation}
The minimum mean square error (MMSE) estimator incorporates channel
statistics to minimise the expected squared error:
\begin{equation}
  \hat{\bH}^{\mathrm{MMSE}} = \bR_H \!\left(\bR_H + \frac{\sigma_n^2}{\rho}\bI\right)^{\!-1}\!\hat{\bH}^{\mathrm{LS}}
\end{equation}
where $\bR_H$ is the channel covariance matrix and $\rho$ is the pilot
signal-to-noise ratio (SNR). The MMSE estimator achieves significantly
better performance but requires knowledge of $\bR_H$, which is often
estimated from historical observations~\cite{Li2025survey}.

Both methods, however, produce \textit{instantaneous} estimates that
do not account for the temporal evolution of the channel. They are
estimation methods, not prediction methods, and thus do not address
the channel aging problem.

%% ═══════════════════════════════════════════════════════════════════════
\section{The Channel Aging Problem}

\subsection{Mathematical Formulation}

The time-varying wireless channel is commonly modelled using Jake's
isotropic scattering model~\cite{Kim2021kalman}. Under this model, the
channel at time step $t+\delta$ is related to the channel at time $t$
through the normalised Doppler frequency $f_D = f_c v / c$, where $f_c$
is the carrier frequency, $v$ is the user velocity, and $c$ is the speed
of light. The temporal correlation coefficient is:
\begin{equation}
  \rho(\delta) = J_0\!\left(2\pi f_D \delta T_s\right)
\end{equation}
where $J_0(\cdot)$ is the zeroth-order Bessel function of the first kind
and $T_s$ is the sampling interval. As $\delta$ increases --- as more
time passes between channel measurement and data transmission ---
$\rho(\delta)$ decays, quantifying how quickly the channel becomes
outdated.

The normalised Doppler frequency $f_n = f_D T_s$ is a key parameter.
Table~\ref{tab:doppler} illustrates the channel coherence degradation
at different user velocities for a 3.5\,GHz 5G-NR carrier.

\begin{table}[htbp]
  \centering
  \caption{Normalised Doppler frequency and correlation coefficient at
  3.5\,GHz for different user velocities (5\,ms processing delay assumed).}
  \label{tab:doppler}
  \renewcommand{\arraystretch}{1.25}
  \begin{tabular}{cccc}
    \toprule
    \textbf{Velocity (km/h)} & \textbf{$f_D$ (Hz)} & \textbf{$f_n$} & \textbf{$\rho(\delta)$} \\
    \midrule
    3   (pedestrian)   & 9.7   & $4.9\times10^{-5}$ & $\approx 1.00$ \\
    30  (urban)        & 97.2  & $4.9\times10^{-4}$ & $\approx 0.98$ \\
    60  (suburban)     & 194.4 & $9.7\times10^{-4}$ & $\approx 0.94$ \\
    120 (highway)      & 388.9 & $1.9\times10^{-3}$ & $\approx 0.77$ \\
    300 (high-speed rail) & 972.2 & $4.9\times10^{-3}$ & $\approx 0.32$ \\
    \bottomrule
  \end{tabular}
\end{table}

\subsection{Impact on Spectral Efficiency}

The work of Zheng \textit{et al.}~\cite{Zheng2021impact} provided the
first comprehensive closed-form analysis of channel aging in cell-free
massive MIMO systems. They derived spectral efficiency (SE) expressions
incorporating both imperfect channel estimation and aging under
spatially correlated channels, finding that:
\begin{itemize}[nosep]
  \item Channel aging causes monotonically decreasing SE as the Doppler
        spread increases.
  \item Cell-free systems, while inherently more resilient to aging
        than small-cell systems due to macrodiversity, still suffer
        meaningful SE losses at moderate-to-high user velocities.
  \item Statistical channel cooperation power control (SCCPC) partially
        mitigates aging effects but its gains erode as aging intensifies.
\end{itemize}

Similarly, the empirical measurement study by
L\"{o}schenbrand \textit{et al.}~\cite{Loschenbrand2022measurement}
provided real-world validation of channel aging effects in a widely
distributed cell-free array with 32 BS antennas and vehicular users
moving at 15--60\,km/h, confirming analytical predictions and
demonstrating that aging is a practical concern, not merely a
theoretical artefact.

\subsection{Cell-Free Specific Aggravating Factors}

In a conventional cellular system, the antenna array is co-located,
and the spatial structure of the channel is well-defined relative to a
single array geometry. In a cell-free system, APs are deployed on lamp
posts, building facades, and bus shelters --- positions determined by
urban infrastructure rather than signal processing optimality. As a
result, the inter-AP spatial correlation structure is \textit{irregular},
and channel prediction models designed for regular antenna array
geometries do not directly transfer~\cite{DeFilippo2024cellfree}.

The work of Freitas \textit{et al.}~\cite{Freitas2025ekf} was among
the first to explicitly address channel aging in cell-free MIMO from a
prediction perspective, proposing an extended Kalman filter (EKF)-based
directional CSI prediction scheme operating under high Doppler
conditions. While effective, the EKF approach requires precise
parametric channel modelling, which is difficult to maintain as the
propagation environment changes dynamically.

%% ═══════════════════════════════════════════════════════════════════════
\section{Cell-Free Massive MIMO: Architecture and Challenges}

\begin{figure}[htbp]
  \centering
  \includegraphics[width=0.90\linewidth]{figures/fig_architecture.pdf}
  \caption{Architecture comparison: (a) conventional cellular network
  with one base station per cell; (b) cell-free network with distributed
  APs cooperating to serve all users simultaneously without cell
  boundaries.}
  \label{fig:architecture}
\end{figure}

\subsection{System Model}

As illustrated in Figure~\ref{fig:architecture}, a cell-free massive
MIMO system consists of $L$ single- or multi-antenna APs distributed
across a geographical area, connected via fronthaul links to a CPU.
All $L$ APs simultaneously serve $K$ user equipment (UE) devices using
the same time-frequency resources. The received signal at the CPU from
user $k$ on subcarrier $m$ is:
\begin{equation}
  \mathbf{y}_{k,m} = \sum_{l=1}^{L} \mathbf{h}_{lk,m}\, x_k
    + \sum_{j \neq k} \sum_{l=1}^{L} \mathbf{h}_{lj,m}\, x_j + \mathbf{n}
\end{equation}
where $\mathbf{h}_{lk,m} \in \CC^{N_l}$ is the channel vector between
AP $l$ and user $k$ on subcarrier $m$, $x_k$ is the transmitted symbol
from user $k$, and $\mathbf{n}$ is additive noise.

A comprehensive survey of this architecture, covering fronthaul design,
scalability, and 5G/6G integration, is provided in
\cite{Elhoushy2022survey}. The seminal work by Ngo \textit{et al.}
\cite{Ngo2017cellfree} demonstrated that cell-free systems offer
5-fold and 10-fold improvements over conventional massive MIMO and
small-cell systems, respectively, in terms of 95th-percentile per-user
throughput. More recent work has extended the architecture to integrate
reconfigurable intelligent surfaces (RIS)~\cite{Shi2024ris} and to
support ultra-reliable low-latency communications (URLLC) for
6G~\cite{Ajmal2024survey}.

\subsection{Channel Estimation in Cell-Free Systems}

Channel estimation in cell-free systems follows a TDD pilot-based
protocol. In the uplink training phase, each user transmits pilot
sequences and each AP estimates the channels from all users
simultaneously. The joint channel matrix across all APs is:
\begin{equation}
  \bH = \begin{bmatrix} \bH_1^\top & \bH_2^\top & \cdots & \bH_L^\top \end{bmatrix}^\top
  \in \CC^{\sum_l N_l \times K}
\end{equation}
where $\bH_l \in \CC^{N_l \times K}$ is the channel sub-matrix for
AP $l$.

A key challenge is \textit{pilot contamination}: when the number of
users $K$ exceeds the pilot sequence length $\tau_p$, pilots must be
reused across users, causing inter-user interference in channel
estimates. While pilot contamination has been extensively studied in
conventional massive MIMO~\cite{Guo2024joint}, it interacts with
channel aging in cell-free systems in ways that are not yet fully
characterised in the literature.

%% ═══════════════════════════════════════════════════════════════════════
\section{Model-Based Channel Prediction Methods}

\subsection{Autoregressive Models}

The autoregressive (AR) model represents the channel as a linear
combination of its past values plus a stochastic innovation:
\begin{equation}
  \bh_t = \sum_{p=1}^{P} \bA_p\, \bh_{t-p} + \mathbf{e}_t
\end{equation}
where $\bA_p$ are the AR coefficient matrices of order $P$ and
$\mathbf{e}_t$ is a zero-mean white noise innovation term. The AR
parameters are estimated from observed channel data using the
Yule-Walker equations or least squares~\cite{Kim2021kalman}.

AR models provide a computationally efficient prediction baseline.
However, their performance degrades rapidly under high mobility because
the channel correlation structure changes faster than the AR parameters
can be updated. Furthermore, they assume stationarity over the
estimation window, which fails in environments with abrupt changes in
the scattering geometry.

\subsection{Kalman Filter-Based Prediction}

The Kalman filter (KF) is the optimal linear predictor for linear
Gaussian dynamical systems. By casting the AR model in state-space
form, the KF recursively updates the channel prediction as new pilot
observations arrive~\cite{Kashyap2017kalman}:
\begin{align}
  \hat{\bh}_{t|t-1} &= \bF\,\hat{\bh}_{t-1|t-1} \\
  \hat{\bh}_{t|t}   &= \hat{\bh}_{t|t-1} + \bK_t\!\left(\bh_t^{\mathrm{obs}} - \hat{\bh}_{t|t-1}\right)
\end{align}
where $\bF$ is the state transition matrix derived from the AR
coefficients, and $\bK_t$ is the Kalman gain.

The work of Kashyap \textit{et al.}~\cite{Kashyap2017kalman}
established the Kalman filter as an effective channel predictor for TDD
massive MIMO, showing that ARMA-based predictors can substantially
increase pilot spacing while maintaining acceptable spectral efficiency.
The vector Kalman filter (VKF) of Kim \textit{et al.}~\cite{Kim2021kalman}
extended this to the multi-antenna case and provided a comparative
benchmark against machine learning approaches.

The EKF-based predictor of Freitas \textit{et al.}~\cite{Freitas2025ekf}
represented the first dedicated Kalman-based predictor for cell-free
MIMO. However, Kalman-based methods share a fundamental limitation:
they require channel model parameters (AR coefficients, noise covariance)
to be estimated and updated in real time, creating a computational
overhead that scales with the number of APs and users --- a significant
scaling problem in cell-free systems with $L = 50$--$200$ APs.

\subsection{Wiener Filter and ARMA Predictors}

Wiener filtering provides the minimum mean-square-error linear predictor
of a stationary stochastic process. While theoretically optimal for
stationary channels, its implementation requires full knowledge of the
channel's power spectral density, which is environment-dependent and
must be estimated adaptively. ARMA-based predictors approximate the
Wiener filter by fitting a higher-order model to the observed temporal
correlation~\cite{Kashyap2017kalman}. These methods show good
performance in low-to-medium mobility scenarios but degrade at high
velocities where the stationarity assumption breaks down.

%% ═══════════════════════════════════════════════════════════════════════
\section{Deep Learning for Channel Estimation and Feedback}

\subsection{CsiNet: The Foundational Deep Learning CSI Architecture}

The breakthrough paper by Wen \textit{et al.}~\cite{Wen2018csinet}
introduced CsiNet, the first deep learning architecture for CSI
compression and feedback in FDD massive MIMO. CsiNet uses a CNN-based
autoencoder: the encoder (at the user side) compresses the channel
matrix into a compact codeword, and the decoder (at the base station)
reconstructs the full CSI from the codeword.

CsiNet demonstrated that CNNs can learn efficient representations of
the spatial structure of CSI matrices, significantly outperforming
compressed sensing-based methods at high compression ratios. Its
success established the paradigm of treating the channel matrix as an
image-like 2D structure amenable to convolutional processing --- a
paradigm directly applicable to channel prediction.

Building on CsiNet, Wang \textit{et al.}~\cite{Wang2019csinetlstm}
introduced CsiNet-LSTM, which augments the CsiNet architecture with
LSTM cells to exploit temporal correlations across consecutive channel
snapshots. CsiNet-LSTM significantly improves reconstruction quality
in time-varying channels by learning both spatial structure (via CNN)
and temporal dynamics (via LSTM), establishing the hybrid
spatial-temporal architecture as a natural design paradigm for
channel-related deep learning tasks.

\subsection{Deep CNN for Wideband Channel Estimation}

Dong \textit{et al.} proposed a deep CNN specifically for wideband
millimetre wave (mmWave) massive MIMO channel estimation, exploiting
the frequency correlation structure across OFDM subcarriers. Their
architecture processes the channel matrix across antenna and frequency
dimensions simultaneously, extracting both spatial and spectral
features. This work demonstrated that the multi-dimensional structure
of the channel matrix
($N_{\mathrm{ant}} \times N_{\mathrm{sub}} \times N_{\mathrm{time}}$)
is naturally suited to 3D convolutional processing.

A lightweight variant was later proposed by Tian
\textit{et al.}~\cite{Tian2025hybrid}, who combined CNN with LSTM
in a regularisation-enhanced hybrid architecture for massive MIMO
channel estimation, demonstrating improved robustness to noise while
maintaining lower computational complexity than Transformer-based
competitors.

\subsection{Joint Channel Estimation and Feedback}

Guo \textit{et al.}~\cite{Guo2024joint} proposed a unified deep
learning framework that jointly performs channel estimation,
compression, and feedback in a single end-to-end trained network,
avoiding the suboptimality introduced by separately optimising each
module. Their framework achieves better spectral efficiency than
two-stage approaches while reducing overall latency --- an important
property for channel-aging-aware systems.

%% ═══════════════════════════════════════════════════════════════════════
\section{Machine Learning for Channel Prediction}

\subsection{Early ML-Based Prediction: MLP and Basic CNN}

The systematic comparison study of Kim \textit{et al.}~\cite{Kim2021kalman}
was the first to compare Kalman filter and machine learning approaches
for massive MIMO channel prediction on realistic 3GPP spatial channel
model (SCM) channels. Their results established several important benchmarks:
\begin{itemize}[nosep]
  \item A multilayer perceptron (MLP) with LMMSE pre-processing
        approaches the performance of the VKF at low complexity.
  \item A CNN-AR architecture (CNN with autoregressive pre-processing)
        outperforms CNN-RNN in low-to-medium mobility scenarios.
  \item All ML-based architectures degrade at high normalised Doppler
        ($f_n > 0.2$), indicating the need for more powerful models.
\end{itemize}

\begin{figure}[htbp]
  \centering
  \includegraphics[width=0.90\linewidth]{figures/fig_pipeline.pdf}
  \caption{Proposed CNN-based channel prediction pipeline: a sequence of
  past CSI matrices is fed into the CNN, which outputs the predicted
  future channel state $\hat{\bH}_{t+k}$ for use in precoding.}
  \label{fig:pipeline}
\end{figure}

\subsection{CNN-Based Real-Time Prediction for High-Speed Scenarios}

Targeting high-speed railway and vehicular scenarios, a novel real-time
CNN channel prediction algorithm was proposed in~\cite{CNN2022realtime},
which uses reference signal-based online training to continuously adapt
to non-stationary channel variations. The algorithm updates its model
using only the most recent pilot observations, enabling tracking of
fast time-varying channels in LTE and 5G OFDM systems without
requiring pre-collected training datasets.

This work demonstrated that CNNs can learn the temporal feature
structure of channels directly from reference signals, making them
suitable for online deployment --- a property important for cell-free
systems where channel conditions vary continuously across the distributed
AP topology.

\subsection{Wideband Channel Prediction and Frequency Correlation}

The machine learning channel prediction study of Ko
\textit{et al.}~\cite{Ko2022wideband} explored wideband massive MIMO
channel prediction, proposing domain transformation techniques to reduce
redundancy across subcarrier channels. Their joint learning approach
trains a single predictor across all subcarriers, exploiting frequency
correlation to reduce data requirements and improve generalisation. This
work also identified ML-based channel prediction for cell-free massive
MIMO as an explicit future research direction, noting that the
combination of wideband operation, distributed AP topology, and high
mobility creates a prediction problem qualitatively different from
single-cell massive MIMO.

\subsection{Hybrid CNN-LSTM Architectures}

The work of Nguyen \textit{et al.}~\cite{Nguyen2023cnnlstm} proposed
a CNN-LSTM model for channel estimation in RIS-NOMA assisted 6G
networks, demonstrating that combining CNN (for spatial feature
extraction) with LSTM (for temporal sequence modelling) yields superior
performance compared to either component alone. The model was shown to
be robust to varying system parameters including transmit SNR, power
allocation, and the number of RIS elements --- a form of generalisation
desirable for cell-free systems where channel statistics vary across
AP locations.

Building on similar principles, the uplink OFDM channel predictor
of De Filippo \textit{et al.}~\cite{DeFilippo2024cellfree} extended
the CNN-LSTM encoder-decoder framework with skip connections for 6G
non-terrestrial networks, predicting the channel frequency response
matrix for entire OFDM slots without pilots. While targeted at
satellite scenarios, this work directly demonstrated the viability of
CNN-LSTM prediction for distributed, multi-AP scenarios.

\subsection{Mobile Cell-Free MIMO: Prediction Challenges}

The mobile cell-free massive MIMO challenges survey by Zheng
\textit{et al.}~\cite{Zheng2023mobile} reviewed prediction methods
in the context of distributed systems. They noted that both Kalman
filter and ML-based predictors degrade in high-mobility scenarios
beyond prediction horizons of approximately 0.1--0.3 wavelengths.
They proposed the concept of the ``predictor antenna'' --- a
physically forward-deployed antenna that observes the channel the
mobile unit will encounter in the near future --- as an alternative
paradigm, but acknowledged that this approach requires hardware
modifications to the AP infrastructure, limiting its practical appeal.

%% ═══════════════════════════════════════════════════════════════════════
\section{Transformer and Attention-Based Approaches}

\subsection{Transformer-Based Channel Prediction}

The Transformer-based channel predictor of Jiang \textit{et al.}
\cite{Jiang2022transformer} represented a significant advance,
proposing a parallel synchronous prediction framework that simultaneously
predicts multiple future channel frames rather than sequentially
predicting one step at a time. Their hybrid attention mechanism combines
self-attention across the time dimension with positional encoding
adapted to channel statistics, achieving prediction NMSE below
$-15$\,dB for user speeds up to 60\,km/h on 3GPP CDL-B and UMa
channel models.

However, the Transformer architecture introduced several practical
limitations in the cell-free context:
\begin{enumerate}[nosep]
  \item Computational complexity is $\mathcal{O}(L^2)$ in sequence
        length, which becomes prohibitive for long prediction horizons.
  \item The model size (millions of parameters) is incompatible with
        deployment on resource-constrained edge hardware at distributed APs.
  \item The model requires extensive matched-condition training data
        and exhibits reduced generalisation across different channel
        environments.
\end{enumerate}

A subsequent study on space-time-frequency correlation exploitation
for cell-free MIMO~\cite{Ko2022wideband} reinforced this conclusion,
noting that FreqConv and SpaceConv based Transformer encoders --- while
achieving excellent accuracy in centralised systems --- face practical
difficulty in extracting spatial correlations under the irregular AP
deployment geometry characteristic of cell-free networks.

\subsection{Attention Mechanisms for Selective Estimation}

Gao \textit{et al.}~\cite{Gao2021attention} proposed an attention-aided
deep learning framework for massive MIMO channel estimation that uses a
``divide-and-conquer'' strategy: an attention mechanism first identifies
the most informative angular regions of the channel, then focuses
estimation resources on those regions. This attention-guided approach
substantially improves estimation accuracy for channels with narrow
angular spread at modest complexity overhead. The result suggests that
selective attention mechanisms can complement lightweight CNN
architectures without prohibitive computational cost --- an idea
potentially applicable in the proposed cell-free prediction design.

%% ═══════════════════════════════════════════════════════════════════════
\section{AI-Native Physical Layer Design for 6G}

The concept of AI-native physical layer design represents a paradigm
shift in wireless communication: rather than applying machine learning
as a post-hoc optimisation on conventionally designed systems, the
physical layer is designed from the ground up to be learned and adapted
by AI models. This vision, articulated in the 6G research community
and partly incorporated into 3GPP Release 18 standardisation activities
for 5G-Advanced, encompasses AI-based channel estimation, AI-based
channel prediction, AI-based beamforming, and AI-based modulation and
coding~\cite{overview6G2024}.

In this context, lightweight CNN-based channel prediction is not merely
a technical optimisation problem but a fundamental building block of
the AI-native physical layer. The proposed research contributes to
this vision by demonstrating that computationally efficient deep
learning can replace or substantially outperform classical model-based
predictors even in the most demanding deployment scenario: the
spatially distributed, high-mobility cell-free network.

The challenges and opportunities of machine learning for future wireless
communications, including channel prediction for cell-free systems,
are comprehensively reviewed in~\cite{ML2025future}. That survey
explicitly lists wideband CSI prediction for cell-free massive MIMO
exploiting space-time-frequency correlations as an open research
problem, directly validating the novelty and timeliness of the
proposed project.

The broader 6G AI and communication landscape is surveyed
in~\cite{overview6G2024}, which further contextualises the need for
edge-deployable, lightweight inference models. Given that cell-free
networks push intelligence to the network edge --- to individual APs
with constrained compute resources --- the ``lightweight'' requirement
in the proposed problem statement is not merely a performance goal but
a fundamental deployment constraint.

%% ═══════════════════════════════════════════════════════════════════════
\section{Performance Comparison and Research Gap Analysis}

\subsection{NMSE Comparison Across Methods}

Figure~\ref{fig:nmse} presents a representative comparison of the
normalised mean square error (NMSE) achieved by the major channel
prediction approaches reviewed in this survey under low and high
mobility conditions. Values are drawn from the simulation results
reported in the cited papers and are presented as illustrative rather
than directly comparable, since exact results depend on simulation
parameters, channel models, and system configurations.

\begin{figure}[htbp]
  \centering
  \includegraphics[width=0.88\linewidth]{figures/fig_nmse.pdf}
  \caption{Representative NMSE comparison (lower is better) across
  channel prediction approaches under low and high user mobility.
  Values are approximate, drawn from the cited literature; exact
  figures depend on simulation setup.}
  \label{fig:nmse}
\end{figure}

The key observation is that while Transformer-based models achieve the
best NMSE in centralised massive MIMO, their computational footprint
is incompatible with edge-deployed cell-free APs. The CNN-AR
architecture of Kim \textit{et al.}~\cite{Kim2021kalman} achieves
strong performance at moderate mobility with significantly lower
complexity, making it the most architecturally appropriate starting
point for a lightweight cell-free predictor.

\subsection{Comprehensive Method Comparison}

Table~\ref{tab:comparison} summarises the key properties of the
reviewed prediction methods along dimensions relevant to cell-free
deployment.

\begin{table}[htbp]
  \centering
  \caption{Comparison of channel prediction methods relevant to
  cell-free massive MIMO deployment. Ratings: \textbf{H}=High,
  \textbf{M}=Medium, \textbf{L}=Low, \textbf{Y}=Yes, \textbf{N}=No.}
  \label{tab:comparison}
  \renewcommand{\arraystretch}{1.3}
  \begin{tabular}{
    >{\raggedright\arraybackslash}p{3.3cm}
    >{\centering\arraybackslash}p{1.5cm}
    >{\centering\arraybackslash}p{1.5cm}
    >{\centering\arraybackslash}p{1.6cm}
    >{\centering\arraybackslash}p{1.8cm}
    >{\centering\arraybackslash}p{1.8cm}
  }
    \toprule
    \textbf{Method} &
    \textbf{Pred.\newline Accuracy} &
    \textbf{Complexity} &
    \textbf{Multi-step} &
    \textbf{CF-MIMO\newline tested} &
    \textbf{Edge\newline deployable} \\
    \midrule
    AR Model                & L--M & Very L  & N       & N   & Y \\
    Kalman Filter           & M    & M       & Limited & N   & Part. \\
    EKF~\cite{Freitas2025ekf} & M  & M       & N       & Y   & Part. \\
    MLP                     & M    & L       & N       & N   & Y \\
    CsiNet~\cite{Wen2018csinet} & M--H & M   & N       & N   & Part. \\
    CsiNet-LSTM~\cite{Wang2019csinetlstm} & H & M & Y  & N   & Part. \\
    CNN-AR~\cite{Kim2021kalman} & H & L--M   & Limited & N   & Y \\
    CNN-LSTM~\cite{Nguyen2023cnnlstm} & H & M & Y     & N   & Part. \\
    Transformer~\cite{Jiang2022transformer} & VH & H & Y & N & N \\
    \midrule
    \textbf{Proposed}       & \textbf{H} & \textbf{L} & \textbf{Y} & \textbf{Y} & \textbf{Y} \\
    \bottomrule
  \end{tabular}
\end{table}

\subsection{Identified Research Gap}

The survey reveals a specific and well-bounded research gap:
\textbf{no existing work has proposed a multi-step predictive CNN
architecture specifically designed for the spatially irregular topology
of distributed cell-free massive MIMO systems, with computational
efficiency suitable for edge deployment at each AP.}

Specifically, existing works fail to simultaneously satisfy all of the
following requirements:

\begin{enumerate}
  \item \textbf{Multi-step prediction.} Most CNN and Kalman-based
        works predict only one step ahead. Multi-step prediction is
        critical in cell-free systems where fronthaul delays compound
        the inherent channel aging.

  \item \textbf{Spatially irregular AP topology.} Cell-free APs are
        irregularly placed. Standard CNNs assume regular grid structure.
        A CNN architecture accounting for the irregular geometry of
        distributed APs is required.

  \item \textbf{Computational efficiency.} Transformer and large
        CNN-LSTM models cannot be deployed on resource-constrained AP
        hardware in real time. A lightweight architecture with
        sub-millisecond inference and a parameter count in the tens of
        thousands --- rather than millions --- is needed.

  \item \textbf{High user mobility.} Most ML-based predictors are
        evaluated at low-to-medium mobility ($< 30$\,km/h). High
        mobility scenarios ($> 60$\,km/h), relevant for vehicular 6G
        use cases, remain a significant challenge for lightweight models.
\end{enumerate}

The proposed research directly and completely addresses all four of these
requirements simultaneously, which no prior work has done.

%% ═══════════════════════════════════════════════════════════════════════
\section{Conclusions}

This literature survey has traced the evolution of channel prediction
research from classical AR and Kalman filter methods, through the
landmark CsiNet architecture, to state-of-the-art Transformer-based
predictors. The key findings are summarised as follows:

\begin{enumerate}
  \item \textbf{Channel aging is a critical, practically validated
        problem} in cell-free massive MIMO, with compounding causes
        beyond those in conventional systems.

  \item \textbf{Classical model-based predictors} (AR, Kalman) are
        computationally tractable but rely on stationarity assumptions
        that fail in dynamic environments and at high mobility.

  \item \textbf{Deep learning} has substantially advanced channel
        estimation and feedback, with CsiNet establishing the paradigm
        of CNN-based spatial feature extraction for CSI matrices.

  \item \textbf{Hybrid CNN-LSTM architectures} represent the current
        practical sweet spot, combining spatial feature extraction with
        temporal sequence modelling at moderate computational cost.

  \item \textbf{Transformer-based predictors} achieve the best accuracy
        but are computationally unsuitable for edge deployment at
        distributed cell-free APs.

  \item \textbf{The specific problem of lightweight, multi-step channel
        prediction for spatially irregular, distributed cell-free MIMO
        under high mobility remains an open research problem,} as
        confirmed by multiple 2023--2025 surveys.
\end{enumerate}

The proposed research --- a lightweight, computationally efficient CNN
designed for the spatially irregular topology of cell-free massive MIMO,
performing multi-step forward channel state prediction under high user
mobility --- directly addresses this gap and makes a timely,
well-motivated contribution to the 6G research agenda.

%% ═══════════════════════════════════════════════════════════════════════
\newpage
\bibliographystyle{IEEEtran}
\bibliography{references}

\end{document}
